\section{Envelope constructions}
\label{sec3:constructions}

Each of the three envelope construction methods follows the same two stage calibration procedure: the shape-discovery set $S_1$ is 
used to determine the geometric shape of the envelope, then the size scaling set $S_2$ is used to determine the scaling factor 
$\hat t$. Each construction assigns a scalar score, $\tau(v,m)$, to a score vector $v$ with observed set $m$. A candidate class is 
then included in the prediction set when $\tau(v,m)\le\hat t$, where $\hat t$ is calculated on $S_2$ as in Section~\ref{sec:setup}. 
For both the collapsed and radial methods, we define $\tau(v)=\inf\{t>0: v\in t\cdot E\}$ using an envelope $E$ which is star shaped 
with respect to the origin and give a closed form expression for it. The strip method works in the same way, except that the 
single scaling factor $\hat t$ multiplies the limits of the staircase rather than dilating the whole region, so $\tau$ is defined 
directly as a maximum of ratios.

We first fix a class $h$ and drop the index $h$. $K$ denotes the number of coordinates that remain after removing the coordinates 
that are missing for every sample of $S_1$. For a vector $v\in[0,\infty)^K$ and observed set $m\subseteq\{1,\dots,K\}$, we write 
$v_j$, $j\in m$, for just its observed entries. We use $\varepsilon=10^{-12}$ to prevent division by zero, and $Q_p(\cdot)$ denotes 
the empirical $p$-quantile.


%%%%%%%%%%%%%%%%%%%%%%%%%%%%%%%%%%%%%%%%%%%%%%%%%%%%%%%%%%%%%%%%%%%%%%%%%%%%%%%%%%%%%%%%%%%%%%%%%%%%%%%%%%%%%%%%%%%%%%%%%%%%%%%
\paragraph{Collapsed envelope.} The basic envelope construction simply collapses the $K$-dimensional score vector to a scalar score 
by taking the mean over the observed coordinates, and calibrates a single quantile on $S_1$:
\begin{equation}\label{eq:collapsed}
\bar v=\frac{1}{|m|}\sum_{j\in m}v_j,\qquad \tilde q=Q_{1-\alpha}\bigl(\{\bar v:v\in S_1\}\bigr),\qquad
\tau(v,m)=\frac{\bar v}{\tilde q},\qquad
E=\{v:\bar v\le\tilde q\}.
\end{equation}

This method discards the geometric information that the $K$ coordinates provide. Two vectors with the same mean are treated 
identically whatever the distribution of their coordinates is, so a single coordinate with strong evidence (a very small $v_j$) is 
averaged together with all the others. The accepted region is a halfspace in the observed coordinates. It serves as a mean score baseline 
for the multidimensional constructions below. After calibration on $S_2$, a vector is accepted when $\bar v/\tilde q\le\hat t$, 
which is the same as $\bar v\le\hat t\,\tilde q$. Here $\hat t$ is the $k$-th smallest value of $\bar v/\tilde q$ over $v\in S_2$, 
so $\hat t\,\tilde q$ is just the $k$-th smallest value of of $\bar v$ over $v\in S_2$. 
% This means that $\tilde q$ does not change the prediction set, and for this method $S_1$ is only 
% used to decide which coordinates are kept.



%%%%%%%%%%%%%%%%%%%%%%%%%%%%%%%%%%%%%%%%%%%%%%%%%%%%%%%%%%%%%%%%%%%%%%%%%%%%%%%%%%%%%%%%%%%%%%%%%%%%%%%%%%%%%%%%%%%%%%%%%%%%%%%
 
\paragraph{Radial envelope.} The radial envelope method retains the directional information of the scores. It works by calibrating 
a boundary radius as a function of direction on the unit sphere, which is restricted to the positive orthant since all nonconformity 
scores are nonnegative. We sample $M$ directions $\{u_r\}_{r=1}^M$ uniformly from the positive orthant by drawing independent vectors $g_r\sim|\mathcal 
N(0,I_K)|$ and normalizing them, $u_r=g_r/\|g_r\|$. For each vector $v\in S_1$, we compute its norm $\Vert{}v\Vert{}$ and  
unit direction $d_v = v / \Vert{}v\Vert{}$ using only the observed coordinates.
For each direction $u_r$, we collect the $n_r$ points of $S_1$ within an angular sector of half width $\delta=30^\circ$, that is, 
those with $d_v\cdot u_r\ge\cos\delta$. The directional radius $\tilde q_r$ is then calibrated as
\begin{equation}\label{eq:radii}
\tilde q_r=\begin{cases}
Q_{1-\alpha}\bigl(\{\|v\|:d_v\cdot u_r\ge\cos\delta\}\bigr) & \text{if } n_r\ge n_{\min},\\[1ex]
w_r\,Q_{1-\alpha}\bigl(\{\|v\|:d_v\cdot u_r\ge\cos\delta\}\bigr)+(1-w_r)\,\tilde q_{\mathrm{marg}} & \text{if } 
0<n_r<n_{\min},\\[1ex]
\tilde q_{\mathrm{marg}} & \text{if } n_r=0,
\end{cases}
\end{equation}

where $n_{\min}=\max\{10,\lfloor 0.2\,|S_1|\rfloor\}$ is the minimum number of points required, $w_r=n_r/n_{\min}$, and 
$\tilde q_{\mathrm{marg}}=Q_{1-\alpha}(\{\|s\|:s\in S_1\})$ is the marginal fallback radius. This helps us avoid a hard 
cutoff between sufficient and insufficient local support. To get a smooth boundary, we compute the radius in an arbitrary direction 
$d$ as a weighted average of all $M$ calibrated radii,
\begin{equation}\label{eq:smooth}
R(d)=\sum_{r=1}^{M}\omega_r(d)\,\tilde q_r,\qquad
\omega_r(d)=\frac{\exp(\kappa\,d\cdot u_r)}{\sum_{r'=1}^{M}\exp(\kappa\,d\cdot u_{r'})},
\end{equation}
where $\kappa>0$ controls the angular smoothness (a larger $\kappa$ gets closer to nearest-direction assignment). We evaluate $R$ 
at the directions $a_p\in\{u_1,\dots,u_M,e_1,\dots,e_K\}$, where the $e_j$ are the coordinate axes. This gives us the boundary points 
$b_p=R(a_p)\,a_p$. The envelope $E$ is the union of the boxes $\{v:0\le v\le b_p\}$ over all $p$, and the scaling function has the 
closed form:
\begin{equation}\label{eq:radial}
\tau(v,m)=\min_{p}\;\max_{j\in m}\;\frac{v_j}{b_{pj}+\varepsilon},
\end{equation}
which is $\inf\{t>0:v\in t\cdot E\}$ computed on the observed coordinates. We use $M=250$ and $\kappa=8$ by default.

%%%%%%%%%%%%%%%%%%%%%%%%%%%%%%%%%%%%%%%%%%%%%%%%%%%%%%%%%%%%%%%%%%%%%%%%%%%%%%%%%%%%%%%%%%%%%%%%%%%%%%%%%%%%%%%%%%%%%%%%%%%%%%%

 → 
 →  (twice, in the limits equation)

\paragraph{Strip envelope.} The strip envelope works by defining conditional bounds for each target coordinate $i$ based on the 
bin occupied by another coordinate $j$. We partition the range of the observed $S_1$ values of coordinate $j$ into $NB$ equal width 
bins, and write $\beta_j(x)$ for the bin that contains $x$ (a value below the smallest or above the largest observed $S_1$ value of 
coordinate $j$ is put in the first or last bin, respectively). We work with the level $\alpha'=\alpha/K$, a Bonferroni type correction 
for the $K$ coordinates. It only affects the shape of the envelope, since the coverage is controlled by $\hat t$. For each bin $n$ 
containing $n_{j,n}$ points of $S_1$ with both $v_j$ and $v_i$ observed, the upper limit along coordinate $i\neq j$ is calibrated as
\begin{equation}\label{eq:limits}
L_{j,i,n}=\begin{cases}
Q_{1-\alpha'}\bigl(\{v_i:v\in S_1,\ \beta_j(v_j)=n\}\bigr) & \text{if } n_{j,n}\ge n_0,\\[1ex]
w_{j,n}\,Q_{1-\alpha'}\bigl(\{v_i:v\in S_1,\ \beta_j(v_j)=n\}\bigr)+(1-w_{j,n})\,\tilde q_i & \text{if } 0<n_{j,n}<n_0,\\[1ex]
\tilde q_i & \text{if } n_{j,n}=0,
\end{cases}
\end{equation}
where $\tilde q_i$ is the marginal $(1-\alpha')$ quantile of coordinate $i$ over $S_1$, $n_0$ is the minimum number of points required 
in a bin, and $w_{j,n}=n_{j,n}/n_0$ shrinks the limit of underpopulated bins towards it, which is similar to the blending rule of the 
radial envelope. Then there is a monotonicity pass that makes sure that the limits are nonincreasing in the bin index,
\begin{equation}\label{eq:monotone}
L_{j,i,n}\leftarrow\max_{n'\ge n}L_{j,i,n'}.
\end{equation}
This ensures that a vector with a smaller value (more confident) in the conditioning coordinate $j$ never receives a tighter bound 
on coordinate $i$ than a vector with a larger value in $j$. The envelope $E$ is a coordinatewise monotone staircase, which is the intersection of these 
pairwise constraints. If a coordinate is unobserved, we ignore all the pairwise constraints involving it, so it is neither used as 
a target nor as a conditioning coordinate. The scaling function is
\begin{equation}\label{eq:strip}
\tau(v,m)=\begin{cases}
\displaystyle\max_{\substack{j,i\in m\\ j\ne i}}\ \dfrac{v_i}{L_{j,i,\beta_j(v_j)}+\varepsilon}, & |m|\ge 2,\\[3ex]
\dfrac{v_j}{\tilde q_j+\varepsilon}, & m=\{j\},
\end{cases}
\end{equation}
so a vector is accepted when $v_i\le\hat t\,L_{j,i,\beta_j(v_j)}$ for every pair of observed coordinates. The single scalar 
$\hat t$ multiplies all the limits, while the bins stay fixed. We use $NB=8$ and $n_0=3$ by default.

%%%%%%%%%%%%%%%%%%%%%%%%%%%%%%%%%%%%%%%%%%%%%%%%%%%%%%%%%%%%%%%%%%%%%%%%%%%%%%%%%%%%%%%%%%%%%%%%%%%%%%%%%%%%%%%%%%%%%%%%%%%%%%%

\paragraph{Properties and missing scores.} The three scaling functions share one structural property.

\begin{proposition}\label{prop:monotone}
Fix an observed set $m$ and let $v,v'$ be score vectors with observed set $m$ such that $v'_j\le v_j$ for all $j\in m$. Then 
$\tau(v',m)\le\tau(v,m)$ for each of the three constructions. In particular, every accepted region $\{\tau\le\hat t\}$ is downward 
closed, so if a candidate class is accepted, it stays accepted when its scores improve.
\end{proposition}

\begin{proof}
For \eqref{eq:collapsed} this is clear, because $\tilde q$ does not depend on $v$ and the mean is nondecreasing in each coordinate. 
In \eqref{eq:radial} each ratio $v_j/(b_{pj}+\varepsilon)$ is nondecreasing in $v_j$, and a maximum and a minimum of nondecreasing 
functions is nondecreasing. For \eqref{eq:strip}, $v'_j\le v_j$ gives $\beta_j(v'_j)\le\beta_j(v_j)$ because $\beta_j$ is 
nondecreasing, so $L_{j,i,\beta_j(v'_j)} \ge L_{j,i,\beta_j(v_j)}$ by \eqref{eq:monotone}. Together with $v'_i\le v_i$, every ratio 
in \eqref{eq:strip} gets smaller. The case $|m|=1$ is again a ratio that is nondecreasing in $v_j$.
\end{proof}

The proposition is not needed for the coverage statement of Section~\ref{sec:setup}. For fixed $S_1$ (and fixed directions $u_r$, 
which are drawn independently of $S_2$ and of the new observation), each $\tau$ is a fixed function of $(v,m)$, so the same rank 
argument applies to all three. Monotonicity is only a sanity property of the learned shapes.

For a vector with observed set $m$, the collapsed method averages over $m$, and the radial and strip methods only use the 
coordinates in $m$. If no coordinate is observed, the collapsed method rejects the candidate class, while the radial and strip 
methods accept it. During fitting we require every vector in $S_1$ and $S_2$ to have at least one observed coordinate, so the 
calibration does not cover this case and we make no coverage claim for it.